\documentclass[aspectratio=169]{beamer}
\usepackage[utf8]{inputenc}

% =====================================================================
% THEME & COLORS
% =====================================================================
\usetheme{Madrid}
\usecolortheme{default}
\usepackage{xcolor}
\usepackage{booktabs}
\usepackage{array}
\usepackage{multirow}
\usepackage{tikz}
\usetikzlibrary{shapes.geometric, arrows, positioning, fit, backgrounds}
\usepackage{listings}
% fontawesome5 not used
\usepackage{tcolorbox}
\tcbuselibrary{skins, breakable}
\usepackage{graphicx}
\usepackage{amsmath}
\usepackage{amssymb}

% ---- Custom Colors ----
\definecolor{neo4jGreen}{RGB}{0, 142, 81}
\definecolor{neo4jDark}{RGB}{0, 80, 50}
\definecolor{cassBlue}{RGB}{20, 100, 200}
\definecolor{cassDark}{RGB}{10, 50, 120}
\definecolor{accentOrange}{RGB}{230, 100, 20}
\definecolor{lightGray}{RGB}{245, 245, 245}
\definecolor{codeBackground}{RGB}{30, 30, 30}
\definecolor{codeGreen}{RGB}{80, 200, 120}
\definecolor{codeYellow}{RGB}{255, 220, 80}
\definecolor{codePurple}{RGB}{180, 120, 255}

% ---- Beamer color overrides ----
\setbeamercolor{structure}{fg=neo4jDark}
\setbeamercolor{frametitle}{bg=neo4jDark, fg=white}
\setbeamercolor{title}{fg=white}
\setbeamercolor{subtitle}{fg=neo4jGreen!60!white}
\setbeamercolor{block title}{bg=neo4jDark, fg=white}
\setbeamercolor{block body}{bg=lightGray}
\setbeamercolor{alerted text}{fg=accentOrange}
\setbeamercolor{example text}{fg=neo4jGreen}

% ---- Custom tcolorbox styles ----
\tcbset{
  neo4jbox/.style={
    colback=neo4jGreen!8!white,
    colframe=neo4jGreen!70!black,
    fonttitle=\bfseries\small,
    sharp corners=south,
    boxrule=0.8pt,
    left=4pt, right=4pt, top=3pt, bottom=3pt,
  },
  cassbox/.style={
    colback=cassBlue!8!white,
    colframe=cassBlue!70!black,
    fonttitle=\bfseries\small,
    sharp corners=south,
    boxrule=0.8pt,
    left=4pt, right=4pt, top=3pt, bottom=3pt,
  },
  deepdivebox/.style={
    colback=accentOrange!8!white,
    colframe=accentOrange!80!black,
    fonttitle=\bfseries\small,
    sharp corners,
    boxrule=0.8pt,
    left=4pt, right=4pt, top=3pt, bottom=3pt,
  },
  infobox/.style={
    colback=gray!10!white,
    colframe=gray!60!black,
    fonttitle=\bfseries\small,
    sharp corners,
    boxrule=0.6pt,
    left=4pt, right=4pt, top=3pt, bottom=3pt,
  }
}

% ---- Code listing style ----
\lstdefinestyle{cypher}{
  language=SQL,
  backgroundcolor=\color{codeBackground},
  basicstyle=\ttfamily\scriptsize\color{white},
  keywordstyle=\color{codeYellow}\bfseries,
  commentstyle=\color{gray!60},
  stringstyle=\color{codeGreen},
  breaklines=true,
  frame=none,
  xleftmargin=4pt,
  xrightmargin=4pt,
  aboveskip=2pt,
  belowskip=2pt,
  morekeywords={MATCH, RETURN, WHERE, CREATE, DELETE, MERGE, SET, WITH,
                ORDER, BY, LIMIT, OPTIONAL, CALL, YIELD, UNWIND, AS,
                DESC, ASC, NOT, AND, OR, IN, STARTS, ENDS, CONTAINS,
                ALL, ANY, NONE, SINGLE, EXISTS, CASE, WHEN, THEN, ELSE,
                END, FOREACH, ON, DETACH, SKIP, UNION, DISTINCT, IS,
                NULL, TRUE, FALSE, count, sum, avg, min, max, collect,
                shortestPath, allShortestPaths},
}

\lstdefinestyle{cql}{
  language=SQL,
  backgroundcolor=\color{codeBackground},
  basicstyle=\ttfamily\scriptsize\color{white},
  keywordstyle=\color{codePurple}\bfseries,
  commentstyle=\color{gray!60},
  stringstyle=\color{codeGreen},
  breaklines=true,
  frame=none,
  xleftmargin=4pt,
  xrightmargin=4pt,
  aboveskip=2pt,
  belowskip=2pt,
  morekeywords={CREATE, TABLE, KEYSPACE, USE, INSERT, INTO, VALUES,
                SELECT, FROM, WHERE, PRIMARY, KEY, WITH, AND, OR,
                REPLICATION, DURABLE, WRITES, CLUSTERING, ORDER,
                IF, NOT, EXISTS, ALTER, DROP, TRUNCATE, USING,
                TTL, TIMESTAMP, LIMIT, ALLOW, FILTERING, BATCH,
                BEGIN, APPLY, ON, SET, UPDATE, DELETE, uuid, timeuuid,
                text, int, bigint, boolean, float, double, blob,
                list, set, map, frozen, counter, inet, timestamp,
                QUORUM, ONE, ALL, LOCAL_QUORUM, EACH_QUORUM},
}

% =====================================================================
% META
% =====================================================================
\title{Neo4j \& Apache Cassandra}
\subtitle{An Academic Deep Dive:\\Graph Databases vs.\ Wide-Column Stores}
\author{L\^{a}m Kh\'{a}nh Tr\`{i}nh}
\institute{Ho Chi Minh City University of Technology (HCMUT)\\
           Faculty of Computer Science \& Engineering}
\date{\today}

% =====================================================================
% DOCUMENT
% =====================================================================
\begin{document}

% ------------------------------------------------------------------
% TITLE
% ------------------------------------------------------------------
\begin{frame}
  \titlepage
\end{frame}

% ------------------------------------------------------------------
% TABLE OF CONTENTS
% ------------------------------------------------------------------
\begin{frame}{Table of Contents}
  \tableofcontents[hideallsubsections]
\end{frame}

% ==================================================================
% SECTION 1: THE NOSQL LANDSCAPE
% ==================================================================
\section{The NoSQL Landscape}

\begin{frame}{The Rise of NoSQL: Historical Context}
  \begin{columns}[T]
    \column{0.55\textwidth}
      \textbf{The Relational Era (1970s -- 2000s)}
      \begin{itemize}\small
        \item Edgar Codd's relational model dominated for decades
        \item SQL databases: strong consistency, rigid schemas, vertical scaling
        \item The Web 2.0 explosion shattered their assumptions:
        \begin{itemize}\scriptsize
          \item Billions of users generating semi-structured data at high velocity
          \item Social graphs, time-series logs, documents — all ill-fitting
          \item Cost of vertical scaling became prohibitive
        \end{itemize}
      \end{itemize}
      \vspace{0.3cm}
      \textbf{The NoSQL Response (2004 -- present)}
      \begin{itemize}\small
        \item Google Bigtable (2004), Amazon Dynamo (2007) papers ignited the movement
        \item ``Not Only SQL'' — horizontal scale, flexible schemas, specialized data models
        \item Key design insight: \textit{different problems need different storage engines}
      \end{itemize}
    \column{0.42\textwidth}
      \begin{tcolorbox}[infobox, title=Four NoSQL Families]
        \scriptsize
        \begin{tabular}{@{}ll@{}}
          \textbf{Key-Value} & Redis, DynamoDB \\[2pt]
          \textbf{Document} & MongoDB, CouchDB \\[2pt]
          \textbf{Wide-Column} & \textcolor{cassBlue}{\textbf{Cassandra}}, HBase \\[2pt]
          \textbf{Graph} & \textcolor{neo4jGreen}{\textbf{Neo4j}}, JanusGraph \\
        \end{tabular}
        \vspace{0.2cm}
        \hrule
        \vspace{0.2cm}
        Each family excels at a specific \textit{access pattern}. Choosing wrong leads to fundamental architectural pain.
      \end{tcolorbox}
    \end{columns}
\end{frame}

% ------------------------------------------------------------------
\begin{frame}{CAP Theorem: The Fundamental Trade-off}
  \begin{columns}[T]
    \column{0.50\textwidth}
      \begin{tcolorbox}[infobox, title={CAP Theorem (Brewer, 2000)}]
        \scriptsize
        A distributed system can guarantee \textbf{at most two} of the three:
        \vspace{0.2cm}
        \begin{itemize}
          \item[\textbf{C}] \textbf{Consistency} — every read receives the most recent write or an error
          \item[\textbf{A}] \textbf{Availability} — every request receives a (possibly stale) response
          \item[\textbf{P}] \textbf{Partition Tolerance} — system continues despite network partitions
        \end{itemize}
        \vspace{0.2cm}
        \textit{In practice, P is mandatory for any distributed system. The real choice is C vs.\ A during a partition.}
      \end{tcolorbox}
      \vspace{0.3cm}
      \begin{tcolorbox}[neo4jbox, title=Neo4j Positioning: CA / CP]
        \scriptsize
        Prioritizes \textbf{Consistency}. Causal Clustering uses the Raft consensus protocol to guarantee that all writes are committed to a majority of Core Servers before acknowledging success.
      \end{tcolorbox}
    \column{0.47\textwidth}
      \begin{tcolorbox}[cassbox, title=Cassandra Positioning: AP]
        \scriptsize
        Prioritizes \textbf{Availability} and \textbf{Partition Tolerance}. Accepts \textit{Eventual Consistency} by default. During a network partition, Cassandra prefers to serve potentially stale data rather than refuse the request.
        \vspace{0.2cm}
        The key innovation: \textbf{Tunable Consistency} allows operators to dial the CP/AP balance per query using consistency levels (ONE, QUORUM, ALL).
      \end{tcolorbox}
      \vspace{0.2cm}
      \begin{tcolorbox}[deepdivebox, title={Deep Dive: PACELC Extension}]
        \scriptsize
        The PACELC model (Abadi, 2012) refines CAP: \textit{even without} a partition, systems trade off \textbf{Latency} vs.\ \textbf{Consistency}. Cassandra is PA/EL (prioritize availability + low latency); Neo4j is PC/EC (prioritize consistency).
      \end{tcolorbox}
  \end{columns}
\end{frame}

% ------------------------------------------------------------------
\begin{frame}{ACID vs.\ BASE: Competing Consistency Philosophies}
  \begin{columns}[T]
    \column{0.47\textwidth}
      \begin{tcolorbox}[neo4jbox, title=ACID (Neo4j)]
        \small
        \begin{itemize}
          \item \textbf{Atomicity} — transaction is all-or-nothing
          \item \textbf{Consistency} — data always moves from valid state to valid state
          \item \textbf{Isolation} — concurrent transactions do not interfere
          \item \textbf{Durability} — committed data survives failures
        \end{itemize}
        \vspace{0.2cm}
        \scriptsize Neo4j enforces full ACID on every graph transaction, including multi-hop relationship writes.
      \end{tcolorbox}
    \column{0.47\textwidth}
      \begin{tcolorbox}[cassbox, title=BASE (Cassandra)]
        \small
        \begin{itemize}
          \item \textbf{Basically Available} — system is always responsive
          \item \textbf{Soft state} — state may change over time without input
          \item \textbf{Eventually consistent} — all replicas converge to same value \textit{eventually}
        \end{itemize}
        \vspace{0.2cm}
        \scriptsize Cassandra introduced \textit{Lightweight Transactions} (Paxos-based) for conditional writes when strict ordering is required, but at significant performance cost.
      \end{tcolorbox}
  \end{columns}
  \vspace{0.4cm}
  \begin{tcolorbox}[infobox, title=Key Insight]
    \small
    Neither philosophy is universally superior. ACID is essential for financial transactions; BASE enables global-scale write throughput. The art is \textbf{choosing the right model for the problem}.
  \end{tcolorbox}
\end{frame}

% ------------------------------------------------------------------
\begin{frame}{Data Models: A Taxonomy}
  \small
  \begin{center}
    \renewcommand{\arraystretch}{1.4}
    \begin{tabular}{>{\bfseries}p{2.8cm} p{3.8cm} p{4.0cm}}
      \toprule
      \textbf{Dimension} & \textbf{Graph (Neo4j)} & \textbf{Wide-Column (Cassandra)} \\
      \midrule
      Core abstraction & Nodes \& typed Relationships & Rows with dynamic column families \\
      Schema & Flexible (schema-optional) & Per-table schema, strictly enforced \\
      Primary query & Pattern matching, traversal & Partition key lookup \\
      Joins & Native, $O(1)$ per hop & Forbidden (application-side only) \\
      Relationships & First-class citizens & Encoded as denormalized columns \\
      Ordering & Graph topology & Clustering column defines disk order \\
      \bottomrule
    \end{tabular}
  \end{center}
\end{frame}

% ==================================================================
% SECTION 2: NEO4J — GRAPH DATABASE
% ==================================================================
\section{Neo4j: The Graph Database}

\begin{frame}{Neo4j: History \& Origins}
  \begin{columns}[T]
    \column{0.58\textwidth}
      \begin{itemize}\small
        \item Founded in \textbf{2007} by Emil Eifrem, Johan Svensson, and Peter Neubauer in Malmö, Sweden — initially as a backend for a content management system
        \item Open-sourced under GPLv3 in 2010; community edition remains free
        \item Introduced the \textbf{Cypher query language} (2011), later donated to the openCypher project (2015) for cross-vendor standardization
        \item Version \textbf{4.x} introduced multi-database and Fabric sharding; \textbf{5.x} unified clustering under the Raft-based Causal Clustering model
        \item Currently the \textbf{\#1 ranked graph database} on DB-Engines and the de-facto standard for property graph systems
      \end{itemize}
    \column{0.38\textwidth}
      \begin{tcolorbox}[neo4jbox, title=Key Milestones]
        \scriptsize
        \begin{tabular}{@{}rl@{}}
          \textbf{2007} & Company founded \\
          \textbf{2010} & First stable release \\
          \textbf{2011} & Cypher language introduced \\
          \textbf{2013} & Enterprise clustering \\
          \textbf{2015} & openCypher project \\
          \textbf{2017} & Causal Clustering (Raft) \\
          \textbf{2020} & Neo4j 4.0 (multi-DB) \\
          \textbf{2022} & AuraDB (managed cloud) \\
          \textbf{2023} & Neo4j 5.x (unified arch.) \\
        \end{tabular}
      \end{tcolorbox}
  \end{columns}
\end{frame}

% ------------------------------------------------------------------
\begin{frame}{The Property Graph Model}
  \begin{columns}[T]
    \column{0.50\textwidth}
      \textbf{Four primitive elements:}
      \vspace{0.2cm}
      \begin{tcolorbox}[neo4jbox, title=1. Nodes (Vertices)]
        \scriptsize
        Represent \textbf{entities} (nouns). Can carry one or more \textit{Labels} (e.g., \texttt{:Student}, \texttt{:Course}) and arbitrary key-value \textit{Properties} (e.g., \texttt{name: "Trinh"}, \texttt{gpa: 3.8}).
      \end{tcolorbox}
      \begin{tcolorbox}[neo4jbox, title=2. Relationships (Edges)]
        \scriptsize
        Represent \textbf{connections} (verbs). Always \textit{directed} (but can be traversed both ways). Have exactly one \textit{Type} (e.g., \texttt{:ENROLLED\_IN}) and can also carry \textit{Properties} (e.g., \texttt{since: 2023}).
      \end{tcolorbox}
      \begin{tcolorbox}[neo4jbox, title=3. Labels \& 4. Properties]
        \scriptsize
        \textbf{Labels}: categorization mechanism; one node can have many labels. \textbf{Properties}: schema-optional key-value pairs stored directly on nodes/edges — no separate tables needed.
      \end{tcolorbox}
    \column{0.47\textwidth}
      \begin{tcolorbox}[infobox, title=Why Graph Matters for Relationships]
        \scriptsize
        In a relational database, representing:
        \begin{center}\texttt{Alice $\to$ FRIEND\_WITH $\to$ Bob}\end{center}
        requires a \texttt{friendships} JOIN table and an expensive index scan. Adding a \textit{property} to that friendship (e.g., \texttt{since: 2020}) requires altering the JOIN table. 
        \vspace{0.2cm}
        In Neo4j, relationships are first-class objects stored with their own properties. The relationship record directly embeds pointers to both endpoint nodes — no index needed.
      \end{tcolorbox}
      \vspace{0.3cm}
      \begin{alertblock}{Key Differentiator}
        \small Relationships in Neo4j are \textbf{not derived at query time}. They are \textit{pre-materialized on disk}, making multi-hop traversals independent of total data set size.
      \end{alertblock}
  \end{columns}
\end{frame}

% ------------------------------------------------------------------
\begin{frame}{Index-Free Adjacency: Neo4j's Performance Secret}
  \begin{columns}[T]
    \column{0.52\textwidth}
      \textbf{How traditional databases join:}
      \begin{itemize}\small
        \item SQL uses \textbf{B-Tree indexes} over foreign key columns
        \item To resolve a JOIN: hash/merge probe on the index — $O(\log N)$ per lookup
        \item Each additional hop in a multi-join query multiplies cost
        \item For a 6-hop social graph traversal on 1 billion nodes: \textbf{infeasible}
      \end{itemize}
      \vspace{0.3cm}
      \textbf{How Neo4j traverses:}
      \begin{itemize}\small
        \item Each node record contains a \textbf{physical pointer} to its first relationship record
        \item Each relationship record contains pointers to: the start node, end node, and the \textit{next relationship} of each node (a doubly-linked list)
        \item Traversal = \textbf{pointer chasing} — $O(1)$ per hop regardless of graph size
      \end{itemize}
    \column{0.45\textwidth}
      \begin{tcolorbox}[deepdivebox, title={Deep Dive: Storage Record Layout}]
        \scriptsize
        Node record (15 bytes):
        \vspace{0.1cm}\\
        \texttt{[inUse|nextRelId|nextPropId|labels]}
        \vspace{0.2cm}\\
        Relationship record (34 bytes):
        \vspace{0.1cm}\\
        \texttt{[inUse|firstNode|secondNode|}\\
        \texttt{ relType|firstPrevRelId|}\\
        \texttt{ firstNextRelId|secondPrevRelId|}\\
        \texttt{ secondNextRelId|nextPropId]}
        \vspace{0.2cm}\\
        Fixed-size records enable \textbf{direct byte-offset addressing}: finding record \#N is simply a seek to byte $N \times 15$ — no heap scan, no index lookup.
      \end{tcolorbox}
      \vspace{0.2cm}
      \begin{tcolorbox}[neo4jbox]
        \scriptsize
        \textbf{Complexity comparison for a $k$-hop query:}
        \begin{center}
          SQL: $O(k \cdot \log N)$ \\
          Neo4j: $O(k \cdot d)$
        \end{center}
        where $d$ = average degree (usually small). Total graph size $N$ \textit{does not appear}.
      \end{tcolorbox}
  \end{columns}
\end{frame}

% ------------------------------------------------------------------
\begin{frame}{Neo4j Architecture: Deployment Modes}
  \begin{columns}[T]
    \column{0.48\textwidth}
      \begin{tcolorbox}[neo4jbox, title=Single Instance]
        \scriptsize
        Full ACID. Reads and writes on same machine. Suitable for development, analytics workloads, or small production deployments with low HA requirements.
      \end{tcolorbox}
      \vspace{0.2cm}
      \begin{tcolorbox}[neo4jbox, title=Causal Clustering (HA)]
        \scriptsize
        \textbf{Core Servers} (typically 3 or 5): participate in Raft consensus. All writes go through Raft leader; a majority must acknowledge before commit.
        \vspace{0.1cm}\\
        \textbf{Read Replicas}: asynchronously replicate from Core Servers; serve read queries; can lag behind by a small delta.
        \vspace{0.1cm}\\
        \textbf{Causal Consistency}: a client receives a \textit{bookmark} after a write; subsequent reads present that bookmark to guarantee they see at least that write.
      \end{tcolorbox}
    \column{0.48\textwidth}
      \begin{tcolorbox}[deepdivebox, title={Deep Dive: Raft Consensus in Neo4j}]
        \scriptsize
        Raft ensures that at most one leader exists per term:
        \begin{enumerate}
          \item Leader receives a write transaction
          \item Appends to its local log and broadcasts to followers
          \item Transaction is committed once a \textbf{quorum} ($\lfloor n/2 \rfloor + 1$) acknowledges
          \item Leader responds to client; followers apply asynchronously
        \end{enumerate}
        Failure scenario: if the leader fails, followers elect a new one via randomized timeouts. No data is lost if a majority survived.
        \vspace{0.1cm}\\
        With 3 Core Servers: tolerates \textbf{1 failure}.\\
        With 5 Core Servers: tolerates \textbf{2 failures}.
      \end{tcolorbox}
      \vspace{0.2cm}
      \begin{tcolorbox}[infobox, title=Neo4j Fabric (Sharding)]
        \scriptsize
        For extremely large graphs: Neo4j Fabric routes subgraph queries to \textit{shards} (separate Neo4j instances), combining results in a federated query layer. Each shard maintains its own Raft cluster.
      \end{tcolorbox}
  \end{columns}
\end{frame}

% ------------------------------------------------------------------
\begin{frame}{Neo4j Storage Engine \& Transaction Lifecycle}
  \small
  \begin{columns}[T]
    \column{0.50\textwidth}
      \textbf{Storage layers:}
      \begin{enumerate}\small
        \item \textbf{Transaction Log} — every write is first appended to the WAL (Write-Ahead Log) for durability, analogous to PostgreSQL's WAL
        \item \textbf{Node/Relationship/Property Stores} — fixed-size record files on disk (mmap'd into memory for fast access)
        \item \textbf{Label Index Store} — allows efficient lookup of all nodes with a given label
        \item \textbf{Schema Indexes} — B-Tree, Full-Text (Lucene), or Point indexes over node/relationship properties
        \item \textbf{Page Cache} — configurable in-memory cache of store pages; the primary performance lever
      \end{enumerate}
    \column{0.47\textwidth}
      \begin{tcolorbox}[deepdivebox, title={Deep Dive: MVCC \& Locking}]
        \scriptsize
        Neo4j uses a \textbf{pessimistic locking} model (not MVCC like PostgreSQL):
        \begin{itemize}
          \item Shared locks on \textit{read} operations
          \item Exclusive locks on \textit{write} operations (node/relationship/property records)
          \item Deadlock detection via cycle-detection in the lock manager wait graph
          \item Phantom reads are prevented by \textit{label scan locks}
        \end{itemize}
        \vspace{0.1cm}
        Implication: high-concurrency write workloads on the \textit{same} nodes/relationships can produce lock contention. Design the data model to minimize hot nodes.
      \end{tcolorbox}
      \vspace{0.2cm}
      \begin{alertblock}{Performance Tuning}
        \scriptsize Set \texttt{dbms.memory.pagecache.size} to fit your full store in RAM. Traversals that stay in the page cache are \textit{orders of magnitude} faster than those requiring disk I/O.
      \end{alertblock}
  \end{columns}
\end{frame}

% ------------------------------------------------------------------
\begin{frame}{Cypher Query Language: Design Philosophy}
  \begin{columns}[T]
    \column{0.53\textwidth}
      \begin{tcolorbox}[neo4jbox, title=ASCII-Art Syntax]
        \scriptsize
        Cypher's defining insight: \textbf{graph patterns look like drawings}.
        \vspace{0.2cm}
        \begin{itemize}
          \item \texttt{(n)} — a node variable \texttt{n}
          \item \texttt{(n:Label)} — node with a label
          \item \texttt{(n \{prop: val\})} — node with a property filter
          \item \texttt{-[:TYPE]->} — directed relationship
          \item \texttt{-[:TYPE*1..3]->} — variable-length path (1 to 3 hops)
          \item \texttt{(a)--(b)} — undirected relationship
        \end{itemize}
        \vspace{0.2cm}
        These compose into \textbf{path patterns} that read like sentences:\\
        \texttt{(alice)-[:FRIEND]->(bob)-[:WORKS\_AT]->(acme)}
      \end{tcolorbox}
    \column{0.44\textwidth}
      \begin{tcolorbox}[infobox, title=Core Cypher Clauses]
        \scriptsize
        \begin{tabular}{@{}ll@{}}
          \texttt{MATCH} & Pattern search \\
          \texttt{OPTIONAL MATCH} & Left outer join \\
          \texttt{WHERE} & Predicate filtering \\
          \texttt{RETURN} & Project output \\
          \texttt{CREATE} & Insert nodes/rels \\
          \texttt{MERGE} & Upsert \\
          \texttt{SET / REMOVE} & Update properties \\
          \texttt{DELETE / DETACH} & Remove data \\
          \texttt{WITH} & Pipeline stages \\
          \texttt{UNWIND} & Flatten lists \\
          \texttt{CALL} & Invoke procedures \\
          \texttt{FOREACH} & Iterate side-effects \\
        \end{tabular}
      \end{tcolorbox}
      \vspace{0.2cm}
      \begin{tcolorbox}[deepdivebox, title=openCypher Standard]
        \scriptsize
        Since 2015, Cypher's specification is maintained by the openCypher community. GQL (ISO/IEC 39075, ratified 2024) standardizes graph query languages, heavily influenced by Cypher.
      \end{tcolorbox}
  \end{columns}
\end{frame}

% ------------------------------------------------------------------
\begin{frame}[fragile]{Cypher: CRUD Operations}
  \begin{columns}[T]
    \column{0.50\textwidth}
      \textbf{Create nodes and relationships:}
      \begin{lstlisting}[style=cypher]
// Create two students and a course
CREATE (t:Student {id: "SV001", name: "Trinh",
                   gpa: 3.8, faculty: "CS"})
CREATE (k:Student {id: "SV002", name: "Khoa",
                   gpa: 3.5, faculty: "CS"})
CREATE (c:Course  {code: "CO3021",
                   name: "Database Systems"})

// Create typed, property-bearing relationship
CREATE (t)-[:STUDIES_WITH {since: 2022,
            project: "NoSQL Research"}]->(k)
CREATE (t)-[:ENROLLED_IN {semester: "HK2"}]->(c)
CREATE (k)-[:ENROLLED_IN {semester: "HK2"}]->(c)
      \end{lstlisting}
      \vspace{0.2cm}
      \textbf{Merge (upsert):}
      \begin{lstlisting}[style=cypher]
// Create only if does not exist
MERGE (p:Professor {id: "GV001"})
ON CREATE SET p.name = "Dr. Nguyen",
              p.created = timestamp()
ON MATCH  SET p.lastSeen = timestamp()
      \end{lstlisting}
    \column{0.47\textwidth}
      \textbf{Update and delete:}
      \begin{lstlisting}[style=cypher]
// Update a property
MATCH (s:Student {id: "SV001"})
SET s.gpa = 3.9, s.honors = true

// Remove a property
MATCH (s:Student {id: "SV001"})
REMOVE s.honors

// Delete a node and all its relationships
MATCH (s:Student {id: "SV002"})
DETACH DELETE s
      \end{lstlisting}
      \vspace{0.2cm}
      \begin{tcolorbox}[neo4jbox, title=MERGE Semantics]
        \scriptsize
        \texttt{MERGE} is Neo4j's upsert: it matches the full pattern, and if no match exists, creates it atomically. Essential for idempotent data loading pipelines.
      \end{tcolorbox}
  \end{columns}
\end{frame}

% ------------------------------------------------------------------
\begin{frame}[fragile]{Cypher: Advanced Pattern Matching}
  \begin{columns}[T]
    \column{0.50\textwidth}
      \textbf{Variable-length path traversal:}
      \begin{lstlisting}[style=cypher]
// Friends within 3 hops of Trinh
// who share the same faculty
MATCH (me:Student {name: "Trinh"})
MATCH (me)-[:FRIEND*1..3]-(network:Student)
WHERE network.faculty = me.faculty
  AND network <> me
RETURN network.name,
       length(shortestPath(
         (me)-[:FRIEND*]-(network)
       )) AS hops
ORDER BY hops ASC
LIMIT 10
      \end{lstlisting}
      \vspace{0.2cm}
      \textbf{Aggregation and collection:}
      \begin{lstlisting}[style=cypher]
// Top courses by enrollment within network
MATCH (me:Student {name: "Trinh"})
MATCH (me)-[:FRIEND*1..2]-(peer:Student)
MATCH (peer)-[:ENROLLED_IN]->(c:Course)
RETURN c.name,
       count(DISTINCT peer) AS students,
       collect(DISTINCT peer.name) AS names
ORDER BY students DESC LIMIT 5
      \end{lstlisting}
    \column{0.47\textwidth}
      \textbf{Shortest path algorithms:}
      \begin{lstlisting}[style=cypher]
// Shortest friendship path between
// two students
MATCH (a:Student {id: "SV001"}),
      (b:Student {id: "SV099"})
MATCH p = shortestPath(
  (a)-[:FRIEND*]-(b)
)
RETURN [n IN nodes(p) | n.name]
         AS path,
       length(p) AS degrees
      \end{lstlisting}
      \vspace{0.2cm}
      \begin{tcolorbox}[deepdivebox, title={Deep Dive: Query Planning}]
        \scriptsize
        Cypher queries are compiled by the \textbf{Cypher Planner} into a physical execution plan (observable via \texttt{EXPLAIN} / \texttt{PROFILE}). The planner uses statistics (node/relationship counts) and index availability to choose between:\\
        \textbf{NodeIndexScan}, \textbf{NodeByLabelScan}, \textbf{Expand}, \textbf{ShortestPath}, \textbf{NodeHashJoin}, etc.
      \end{tcolorbox}
  \end{columns}
\end{frame}

% ------------------------------------------------------------------
\begin{frame}[fragile]{Cypher: Graph Algorithms \& APOC}
  \begin{columns}[T]
    \column{0.50\textwidth}
      \textbf{Neo4j Graph Data Science (GDS) library:}
      \begin{lstlisting}[style=cypher]
// Project an in-memory graph
CALL gds.graph.project(
  'social-graph',
  'Student',
  'FRIEND'
)

// Run PageRank to find influence
CALL gds.pageRank.stream('social-graph')
YIELD nodeId, score
RETURN gds.util.asNode(nodeId).name AS name,
       score
ORDER BY score DESC LIMIT 10
      \end{lstlisting}
      \begin{lstlisting}[style=cypher]
// Community detection (Louvain)
CALL gds.louvain.stream('social-graph')
YIELD nodeId, communityId
RETURN communityId,
       count(*) AS members
ORDER BY members DESC
      \end{lstlisting}
    \column{0.47\textwidth}
      \begin{tcolorbox}[neo4jbox, title=GDS Algorithm Categories]
        \scriptsize
        \begin{tabular}{@{}ll@{}}
          \textbf{Centrality} & PageRank, Betweenness \\
          \textbf{Community} & Louvain, LPA, WCC \\
          \textbf{Path} & Dijkstra, A*, Yen's k-shortest \\
          \textbf{Similarity} & Node2Vec, cosine, Jaccard \\
          \textbf{Link Prediction} & Adamic-Adar, Common Neighbors \\
          \textbf{Embeddings} & FastRP, GraphSAGE \\
        \end{tabular}
      \end{tcolorbox}
      \vspace{0.2cm}
      \begin{tcolorbox}[infobox, title=APOC (Awesome Procedures on Cypher)]
        \scriptsize
        APOC provides 450+ utility procedures for:
        \begin{itemize}
          \item Data import (CSV, JSON, GraphQL)
          \item Dynamic Cypher execution
          \item Date/time manipulation
          \item Graph refactoring
          \item Batch processing with \texttt{apoc.periodic.iterate}
        \end{itemize}
      \end{tcolorbox}
  \end{columns}
\end{frame}

% ------------------------------------------------------------------
\begin{frame}{Neo4j Indexing Strategies}
  \begin{columns}[T]
    \column{0.50\textwidth}
      \small
      \textbf{Schema Indexes} (on node properties):
      \begin{itemize}\scriptsize
        \item \textbf{B-Tree Index} — range queries, equality, ORDER BY. Default for most property types.
        \item \textbf{Full-Text Index} — powered by Apache Lucene; supports \texttt{CONTAINS}, fuzzy matching, relevance scoring. Ideal for search on string properties.
        \item \textbf{Point Index} — spatial data types; supports distance and bounding-box queries.
        \item \textbf{Composite Index} — across multiple properties; used when WHERE filters multiple fields simultaneously.
        \item \textbf{Token Lookup Index} — auto-created; enables fast \texttt{NodeByLabelScan}.
      \end{itemize}
      \vspace{0.2cm}
      \textbf{Relationship indexes:} from Neo4j 4.3+, indexes can be created on relationship properties as well.
    \column{0.47\textwidth}
      \begin{tcolorbox}[deepdivebox, title={Deep Dive: Index vs.\ Index-Free Adjacency}]
        \scriptsize
        It is important to understand when each mechanism fires:
        \vspace{0.1cm}
        \begin{itemize}
          \item \textbf{Schema indexes} are used as \textit{entry points} — to find the first node(s) in a query (e.g., \texttt{MATCH (s:Student \{id:"SV001"\})})
          \item Once the entry node is found, all subsequent \texttt{-[:REL]->} traversals use \textbf{index-free adjacency} (pointer following)
        \end{itemize}
        \vspace{0.1cm}
        This hybrid design combines the best of both worlds: fast property lookup via B-Tree + fast graph traversal via physical pointers.
      \end{tcolorbox}
      \vspace{0.2cm}
      \begin{alertblock}{Index Hint}
        \scriptsize Use \texttt{USING INDEX s:Student(id)} in Cypher to force the planner to use a specific index when the auto-selected plan is suboptimal.
      \end{alertblock}
  \end{columns}
\end{frame}

% ------------------------------------------------------------------
\begin{frame}{Neo4j Use Cases \& Anti-Patterns}
  \begin{columns}[T]
    \column{0.48\textwidth}
      \begin{tcolorbox}[neo4jbox, title={[OK] Ideal Use Cases}]
        \small
        \begin{itemize}
          \item \textbf{Fraud ring detection} — identifying circular transaction patterns that look legitimate individually but form suspicious rings when traversed as a graph
          \item \textbf{Real-time recommendations} — ``People who viewed X also viewed Y'' traversed in real time across social and purchase graphs
          \item \textbf{Knowledge graphs} — connecting entities across multiple domains (research, enterprise ontologies, semantic web)
          \item \textbf{Access control} — evaluating complex permission inheritance hierarchies (RBAC, ABAC)
          \item \textbf{Network/IT topology} — impact analysis of infrastructure failures through dependency graphs
          \item \textbf{Master Data Management} — relating products, suppliers, regulations across enterprise silos
        \end{itemize}
      \end{tcolorbox}
    \column{0.48\textwidth}
      \begin{tcolorbox}[deepdivebox, title={[!] Anti-Patterns}]
        \small
        \begin{itemize}
          \item \textbf{Bulk data logging} — massive flat record ingestion (metrics, events, access logs) where no relationship traversal occurs; Cassandra or ClickHouse is far better
          \item \textbf{Simple key-value lookups} — adding graph overhead to what is fundamentally a dictionary lookup
          \item \textbf{High-cardinality, schema-rigid data} — if data is perfectly tabular and queries never need joins, a relational database or columnar store outperforms
          \item \textbf{Supernode problem} — nodes with millions of relationships (e.g., a ``celebrity'' node) degrade traversal performance; requires \textit{relationship type} partitioning strategies
        \end{itemize}
      \end{tcolorbox}
  \end{columns}
\end{frame}

% ==================================================================
% SECTION 3: CASSANDRA — WIDE-COLUMN STORE
% ==================================================================
\section{Apache Cassandra: The Wide-Column Store}

\begin{frame}{Cassandra: History \& Origins}
  \begin{columns}[T]
    \column{0.58\textwidth}
      \begin{itemize}\small
        \item Originally developed by \textbf{Facebook} (2007--2008) to power the Inbox Search feature — required handling billions of messages with no downtime
        \item Open-sourced and donated to the \textbf{Apache Software Foundation} in 2009; graduated to Top-Level Project in 2010
        \item Architectural DNA: combines \textbf{Amazon Dynamo's} distributed hash ring (for availability) with \textbf{Google Bigtable's} data model (for structured storage)
        \item Major adopters include Netflix (replicas across 3 AWS regions), Apple ($>$75,000 nodes), Discord, Uber, and Instagram
        \item \textbf{DataStax Enterprise (DSE)} is the commercial distribution with enterprise add-ons; Astra DB is the managed cloud offering
      \end{itemize}
    \column{0.38\textwidth}
      \begin{tcolorbox}[cassbox, title=Key Milestones]
        \scriptsize
        \begin{tabular}{@{}rl@{}}
          \textbf{2007} & Created at Facebook \\
          \textbf{2008} & Deployed for Inbox Search \\
          \textbf{2009} & Open-sourced (Apache) \\
          \textbf{2010} & Apache TLP graduation \\
          \textbf{2012} & CQL 3.0 released \\
          \textbf{2014} & Lightweight Transactions \\
          \textbf{2016} & SASI indexing \\
          \textbf{2020} & Cassandra 4.0 (audit log) \\
          \textbf{2022} & Cassandra 4.1 (pluggable) \\
        \end{tabular}
      \end{tcolorbox}
      \vspace{0.2cm}
      \begin{tcolorbox}[deepdivebox, title=Scale Fact]
        \scriptsize
        Apple runs Cassandra on over \textbf{75,000 nodes}, storing more than \textbf{10 petabytes} of data — the largest known Cassandra deployment.
      \end{tcolorbox}
  \end{columns}
\end{frame}

% ------------------------------------------------------------------
\begin{frame}{The Wide-Column Data Model}
  \begin{columns}[T]
    \column{0.52\textwidth}
      \small
      \textbf{Conceptual hierarchy:}
      \begin{itemize}
        \item \textbf{Keyspace} — top-level container (like a relational database); defines replication strategy
        \item \textbf{Table} — collection of rows; each table has a declared schema
        \item \textbf{Row} — uniquely identified by its \textit{primary key}; can have different non-null columns per row (wide columns)
        \item \textbf{Primary Key anatomy:}
          \begin{itemize}\scriptsize
            \item \textbf{Partition Key} — determines which node(s) own this row via consistent hashing
            \item \textbf{Clustering Columns} — determine the sorted order of rows \textit{within} a partition on disk
            \item \textbf{Regular Columns} — the data payload
          \end{itemize}
      \end{itemize}
    \column{0.45\textwidth}
      \begin{tcolorbox}[deepdivebox, title={Deep Dive: ``Wide Column'' Explained}]
        \scriptsize
        The term refers to Bigtable's original model where each row can have \textit{millions of columns} — not the same as a relational table. In Cassandra, this manifests as clustering columns that create a sorted, sparse matrix structure within each partition.
        \vspace{0.2cm}\\
        A single Cassandra partition can hold up to \textbf{2 billion cells} (column values), making it ideal for time-series data where each event becomes a new row within the same partition.
      \end{tcolorbox}
      \vspace{0.2cm}
      \begin{tcolorbox}[infobox, title=Important Constraint]
        \scriptsize
        Cassandra does \textbf{not} support:
        \begin{itemize}
          \item JOINs between tables
          \item Subqueries
          \item Arbitrary \texttt{WHERE} clauses (must use partition key)
          \item Aggregations across partitions (without ALLOW FILTERING, which is dangerous)
        \end{itemize}
      \end{tcolorbox}
  \end{columns}
\end{frame}

% ------------------------------------------------------------------
\begin{frame}{Cassandra: The Masterless Ring Architecture}
  \begin{columns}[T]
    \column{0.50\textwidth}
      \textbf{Core architectural principles:}
      \begin{itemize}\small
        \item \textbf{Peer-to-peer (P2P)} — no master node; every node is functionally identical
        \item \textbf{Consistent Hashing} — the partition key is hashed to a 64-bit token; the ring is divided into token ranges, each owned by a node
        \item \textbf{Replication Factor (RF)} — each partition is replicated on RF nodes for fault tolerance; typically RF=3
        \item \textbf{Coordinator} — any node can accept client requests and acts as coordinator for that request, routing sub-requests to the correct replicas
        \item \textbf{Virtual Nodes (vnodes)} — each physical node owns multiple token ranges (default 256), enabling automatic data balancing as nodes join/leave
      \end{itemize}
    \column{0.47\textwidth}
      \begin{tcolorbox}[deepdivebox, title={Deep Dive: Consistent Hashing}]
        \scriptsize
        \textbf{Problem solved}: with naïve modular hashing (\texttt{key \% N}), adding/removing one node \textit{N} invalidates $\approx N/(N+1)$ of all keys.
        \vspace{0.1cm}\\
        \textbf{Solution}: map both nodes and keys onto a fixed-size ring. When a node is added, it only takes responsibility for a \textit{portion} of its neighbor's range. On average, only \texttt{K/N} keys are remapped (where K = total keys, N = nodes).
        \vspace{0.1cm}\\
        Virtual nodes extend this: each physical node's load can be tuned by assigning it more/fewer vnodes without moving data.
      \end{tcolorbox}
      \vspace{0.2cm}
      \begin{alertblock}{Zero Single Point of Failure}
        \scriptsize Because any node can coordinate any request, Cassandra has \textbf{no master to fail}. Rolling upgrades, node additions, and decommissions are all done without downtime.
      \end{alertblock}
  \end{columns}
\end{frame}

% ------------------------------------------------------------------
\begin{frame}{Gossip Protocol \& Failure Detection}
  \begin{columns}[T]
    \column{0.52\textwidth}
      \textbf{Gossip protocol mechanics:}
      \begin{itemize}\small
        \item Each node periodically (every 1 second) selects 1--3 random peers and exchanges \textbf{state information} (token ownership, load, schema version, generation counter)
        \item Messages contain not only direct knowledge but also \textit{known knowledge about other nodes} — rumors spread virally
        \item Convergence time: $O(\log N)$ rounds for full cluster awareness
        \item No central directory needed — each node builds a complete cluster view
      \end{itemize}
      \vspace{0.2cm}
      \textbf{Failure detection ($\Phi$ Accrual):}
      \begin{itemize}\small
        \item Cassandra does \textit{not} use binary alive/dead state
        \item Computes a continuous suspicion level $\Phi$ based on recent heartbeat history
        \item As $\Phi$ rises above threshold (default $\Phi = 8$), the node is considered unreachable
      \end{itemize}
    \column{0.45\textwidth}
      \begin{tcolorbox}[deepdivebox, title={Deep Dive: $\Phi$ Accrual Detector}]
        \scriptsize
        The $\Phi$ value is computed from inter-arrival times of heartbeat messages:
        \[\Phi(t) = -\log_{10}\big(1 - F(t - t_{\text{last}})\big)\]
        where $F$ is the CDF of historical arrival intervals (modeled as a Gaussian distribution). As the gap between expected and actual heartbeat grows, $\Phi$ increases toward infinity.
        \vspace{0.1cm}\\
        This adaptive model avoids false positives in high-latency networks (e.g., cross-datacenter links) by adjusting the threshold automatically based on network conditions.
      \end{tcolorbox}
      \vspace{0.2cm}
      \begin{tcolorbox}[cassbox]
        \scriptsize
        \textbf{Key benefit}: Gossip + $\Phi$ detection enables Cassandra to heal gracefully from transient network blips without marking nodes as permanently failed.
      \end{tcolorbox}
  \end{columns}
\end{frame}

% ------------------------------------------------------------------
\begin{frame}{Replication Strategies in Cassandra}
  \small
  \begin{columns}[T]
    \column{0.50\textwidth}
      \begin{tcolorbox}[cassbox, title=SimpleStrategy]
        \scriptsize
        Places replicas on the next $RF-1$ nodes clockwise on the ring, regardless of rack or datacenter. \textbf{Use only for single-datacenter deployments or development}.
      \end{tcolorbox}
      \vspace{0.2cm}
      \begin{tcolorbox}[cassbox, title=NetworkTopologyStrategy (Production)]
        \scriptsize
        Specifies per-datacenter replication factors. Distributes replicas across different racks within each datacenter for maximum fault isolation.
        \vspace{0.1cm}\\
        \textbf{Example}: \texttt{datacenter1: 3, datacenter2: 3} means 6 total replicas across 2 DCs. A failure of one entire DC still allows the other DC to serve all data.
      \end{tcolorbox}
    \column{0.47\textwidth}
      \begin{tcolorbox}[deepdivebox, title={Deep Dive: Multi-DC Write Path}]
        \scriptsize
        With \texttt{LOCAL\_QUORUM} consistency:
        \begin{enumerate}
          \item Client writes to any coordinator in DC1
          \item Coordinator routes write to all RF replicas in DC1 (synchronously)
          \item Coordinator acknowledges to client once local quorum is met
          \item Replication to DC2 happens \textbf{asynchronously} via hints
        \end{enumerate}
        This achieves low-latency writes for the local DC while maintaining eventual consistency with the remote DC — ideal for active-active multi-region deployments.
      \end{tcolorbox}
      \vspace{0.2cm}
      \begin{tcolorbox}[infobox, title=Hinted Handoff]
        \scriptsize
        If a replica node is temporarily unavailable, the coordinator stores a \textit{hint} (a buffer of the missed write). When the replica comes back online, the hint is replayed to bring it up to date.
      \end{tcolorbox}
  \end{columns}
\end{frame}

% ------------------------------------------------------------------
\begin{frame}{Cassandra Storage Engine: The Write Path}
  \begin{columns}[T]
    \column{0.55\textwidth}
      \small
      \textbf{Step-by-step write path:}
      \begin{enumerate}
        \item \textbf{Commit Log} — write is appended to the per-node commit log (WAL) immediately. This is \textit{sequential I/O} on disk — extremely fast. Ensures durability across crashes.
        \item \textbf{Memtable} — write is also written to an in-memory sorted structure (the Memtable). Reads on recently written data are served from here.
        \item \textbf{SSTable flush} — when the Memtable exceeds a threshold, it is \textit{flushed} to disk as a new, immutable SSTable (Sorted String Table). Flush is also sequential.
        \item \textbf{Bloom Filter} — each SSTable has a Bloom filter in memory; used to quickly determine if a key \textit{might} exist in that SSTable, avoiding unnecessary disk seeks.
        \item \textbf{Compaction} — background process that merges and de-duplicates SSTables, removing tombstones and reclaiming disk space.
      \end{enumerate}
    \column{0.42\textwidth}
      \begin{tcolorbox}[deepdivebox, title={Deep Dive: Tombstones}]
        \scriptsize
        Cassandra does \textbf{not} delete data in-place. Deletions write a special record called a \textit{tombstone} with a deletion timestamp. During compaction, rows whose tombstone timestamp exceeds \texttt{gc\_grace\_seconds} (default: 10 days) are actually removed.
        \vspace{0.1cm}\\
        \textbf{Why 10 days?} If a node was offline for 9 days and missed a delete, it must receive the tombstone via repair before the GC window expires. If the window expires first, the deleted data is resurrected (a ``zombie'' row).
        \vspace{0.1cm}\\
        \textbf{Implication}: designs that generate many tombstones (e.g., writing a ``null'' to update a cell) severely degrade read performance. Prefer \texttt{TTL} instead.
      \end{tcolorbox}
  \end{columns}
\end{frame}

% ------------------------------------------------------------------
\begin{frame}{Cassandra Storage Engine: The Read Path}
  \begin{columns}[T]
    \column{0.52\textwidth}
      \small
      \textbf{Step-by-step read path:}
      \begin{enumerate}
        \item Check \textbf{Row Cache} (if enabled and configured)
        \item Check \textbf{Memtable} — most recently written data served instantly
        \item For each SSTable:
        \begin{enumerate}\scriptsize
          \item Check \textbf{Bloom filter} — skip SSTable if key definitely absent ($O(1)$)
          \item Check \textbf{Partition Summary} — memory-resident sparse index of partition offsets
          \item Consult \textbf{Partition Index} — on-disk structure for exact byte offset of key
          \item Seek \textbf{SSTable data file} — read compressed, sorted data blocks
        \end{enumerate}
        \item \textbf{Merge} results from Memtable and all SSTables; use timestamp to determine winner (last-write-wins)
        \item Apply \textbf{tombstones} to filter deleted rows
      \end{enumerate}
    \column{0.45\textwidth}
      \begin{tcolorbox}[deepdivebox, title={Deep Dive: Read Repair}]
        \scriptsize
        When consistency level requires reading from multiple replicas:
        \begin{enumerate}
          \item Coordinator sends read to all required replicas
          \item Compares digest responses
          \item If replicas disagree: fetches full data from all, picks latest timestamp, and \textit{writes the correct value back} to the stale replicas
          \item Returns the correct value to the client
        \end{enumerate}
        This self-healing mechanism is called \textbf{Read Repair}. It runs in the foreground (blocking) at the configured probability, keeping replicas in sync without explicit repair jobs.
      \end{tcolorbox}
      \vspace{0.2cm}
      \begin{alertblock}{Compaction Strategy}
        \scriptsize \textbf{STCS} (default) — good for write-heavy; \textbf{LCS} — good for read-heavy (fewer SSTables); \textbf{TWCS} — ideal for time-series with TTL.
      \end{alertblock}
  \end{columns}
\end{frame}

% ------------------------------------------------------------------
\begin{frame}{Tunable Consistency: The AP/CP Dial}
  \begin{columns}[T]
    \column{0.50\textwidth}
      \small
      \textbf{Consistency level governs how many replicas must respond:}
      \vspace{0.2cm}
      \begin{center}
        \renewcommand{\arraystretch}{1.3}
        \begin{tabular}{>{\ttfamily\scriptsize}p{2.0cm} p{3.6cm}}
          \toprule
          \textbf{Level} & \textbf{Semantics} \\
          \midrule
          ONE & Only 1 replica must respond (fastest, least consistent) \\
          TWO / THREE & 2 or 3 replicas must respond \\
          QUORUM & Majority: $\lfloor RF/2 \rfloor + 1$ replicas \\
          LOCAL\_QUORUM & Quorum within local DC only \\
          EACH\_QUORUM & Quorum in \textit{each} DC \\
          ALL & All replicas must respond (most consistent, least available) \\
          \bottomrule
        \end{tabular}
      \end{center}
    \column{0.47\textwidth}
      \begin{tcolorbox}[deepdivebox, title={Deep Dive: Strong Consistency Formula}]
        \scriptsize
        With RF replicas, to achieve strong (linearizable) consistency, you need:
        \[R + W > RF\]
        where $R$ = read CL replicas, $W$ = write CL replicas.
        \vspace{0.1cm}\\
        \textbf{Example} (RF=3):
        \begin{itemize}
          \item Write: QUORUM (2) + Read: QUORUM (2) = 4 $>$ 3 \checkmark \textit{Strong}
          \item Write: ONE (1) + Read: ONE (1) = 2 $\not>$ 3 \texttimes \textit{Eventual}
          \item Write: ALL (3) + Read: ONE (1) = 4 $>$ 3 \checkmark \textit{Strong (but low write availability)}
        \end{itemize}
        QUORUM + QUORUM is the typical production setting for most applications.
      \end{tcolorbox}
      \vspace{0.2cm}
      \begin{tcolorbox}[cassbox]
        \scriptsize
        \textbf{Key design freedom}: you can set consistency \textit{per query}. A banking query can use \texttt{QUORUM}; a metrics query can use \texttt{ONE}.
      \end{tcolorbox}
  \end{columns}
\end{frame}

% ------------------------------------------------------------------
\begin{frame}{Query-First Data Modeling: The Core Principle}
  \begin{columns}[T]
    \column{0.50\textwidth}
      \begin{alertblock}{Cassandra's Cardinal Rule}
        \small Model your \textbf{tables around your queries}, not around your entities. Every query should be answerable by a single partition scan.
      \end{alertblock}
      \vspace{0.3cm}
      \textbf{The modeling workflow:}
      \begin{enumerate}\small
        \item Enumerate all application queries (``What does the app need to read?'')
        \item For each query: design a table whose partition key equals the query's primary filter
        \item Choose clustering columns to sort data in the order the query needs
        \item Accept data duplication as a \textit{feature}, not a bug
        \item Define the write path: what tables must be updated on each write event?
      \end{enumerate}
    \column{0.47\textwidth}
      \begin{tcolorbox}[deepdivebox, title={Deep Dive: Why No JOINs?}]
        \scriptsize
        A JOIN in a distributed system requires:
        \begin{enumerate}
          \item Broadcast query to all nodes holding relevant partitions
          \item Shuffle large data volumes across the network
          \item Merge results centrally
        \end{enumerate}
        In a cluster with 100 nodes, this is a \textbf{fan-out of 100} network round-trips. Under high load, this becomes a distributed coordination nightmare.
        \vspace{0.1cm}\\
        \textbf{Cassandra's answer}: pre-compute the join at write time. Denormalize the data so each query only touches one partition. JOINs happen in the \textit{application}, at write time, not at read time.
      \end{tcolorbox}
  \end{columns}
\end{frame}

% ------------------------------------------------------------------
\begin{frame}[fragile]{CQL: Keyspace \& Table Design}
  \begin{columns}[T]
    \column{0.50\textwidth}
      \textbf{Creating a Keyspace:}
      \begin{lstlisting}[style=cql]
CREATE KEYSPACE university
  WITH REPLICATION = {
    'class': 'NetworkTopologyStrategy',
    'datacenter1': 3,
    'datacenter2': 2
  }
  AND DURABLE_WRITES = true;

USE university;
      \end{lstlisting}
      \vspace{0.2cm}
      \textbf{Query-Driven Table: ``Get all enrollments for a student'':}
      \begin{lstlisting}[style=cql]
CREATE TABLE enrollments_by_student (
    student_id  uuid,
    enrolled_at timestamp,
    course_code text,
    course_name text,
    semester    text,
    grade       float,
    PRIMARY KEY ((student_id), enrolled_at)
) WITH CLUSTERING ORDER BY (enrolled_at DESC);
      \end{lstlisting}
    \column{0.47\textwidth}
      \textbf{Query-Driven Table: ``Get top students in a course'':}
      \begin{lstlisting}[style=cql]
-- Separate table for the reverse query
CREATE TABLE students_by_course (
    course_code text,
    gpa         float,
    student_id  uuid,
    student_name text,
    semester    text,
    PRIMARY KEY ((course_code), gpa, student_id)
) WITH CLUSTERING ORDER BY (gpa DESC);
      \end{lstlisting}
      \vspace{0.2cm}
      \begin{tcolorbox}[cassbox, title={Two Tables, Two Queries}]
        \scriptsize
        Both tables store the same relationship data (student $\leftrightarrow$ course enrollment). This \textbf{denormalization} is intentional: each table is optimized for exactly one read query. Write operations update both tables atomically using a \texttt{BATCH} statement.
      \end{tcolorbox}
  \end{columns}
\end{frame}

% ------------------------------------------------------------------
\begin{frame}[fragile]{CQL: Advanced Features}
  \begin{columns}[T]
    \column{0.50\textwidth}
      \textbf{TTL (Time-To-Live) for automatic expiry:}
      \begin{lstlisting}[style=cql]
-- Session token expires in 3600 seconds
INSERT INTO session_tokens
  (user_id, token, created_at)
VALUES (uuid(), 'abc123xyz', toTimestamp(now()))
USING TTL 3600;

-- Per-column TTL
UPDATE user_preferences
USING TTL 86400
SET dark_mode = true
WHERE user_id = 550e8400-e29b-41d4-a716-446655440000;
      \end{lstlisting}
      \vspace{0.2cm}
      \textbf{Lightweight Transactions (Paxos):}
      \begin{lstlisting}[style=cql]
-- Conditional insert (compare-and-set)
INSERT INTO accounts (id, email, balance)
VALUES (uuid(), 'trinh@hcmut.edu.vn', 1000)
IF NOT EXISTS;

-- Conditional update
UPDATE accounts SET balance = 900
WHERE id = ?
IF balance = 1000;
      \end{lstlisting}
    \column{0.47\textwidth}
      \textbf{Collections (List, Set, Map):}
      \begin{lstlisting}[style=cql]
CREATE TABLE student_profile (
    id       uuid PRIMARY KEY,
    name     text,
    tags     set<text>,
    scores   list<float>,
    metadata map<text, text>
);

-- Append to list
UPDATE student_profile
SET scores = scores + [3.8]
WHERE id = ?;

-- Add to set
UPDATE student_profile
SET tags = tags + {'honors', 'research'}
WHERE id = ?;
      \end{lstlisting}
      \begin{tcolorbox}[deepdivebox, title={Deep Dive: LWT Cost}]
        \scriptsize
        Lightweight Transactions use a \textbf{4-round-trip Paxos protocol} (Prepare → Promise → Propose → Commit). They are $4\times$--$8\times$ slower than regular writes. Use them only when you truly need conditional atomicity; avoid overuse.
      \end{tcolorbox}
  \end{columns}
\end{frame}

% ------------------------------------------------------------------
\begin{frame}[fragile]{CQL: Time-Series Data Modeling}
  \begin{columns}[T]
    \column{0.50\textwidth}
      \textbf{Sensor data model — classic time-series:}
      \begin{lstlisting}[style=cql]
CREATE TABLE iot_sensor_readings (
    device_id  uuid,
    -- Bucket by day to control partition size
    day        date,
    recorded_at timestamp,
    temperature float,
    humidity    float,
    pressure    float,
    PRIMARY KEY ((device_id, day), recorded_at)
) WITH CLUSTERING ORDER BY (recorded_at DESC)
  AND default_time_to_live = 2592000; -- 30 days TTL

-- Querying the last hour for device ABC
SELECT recorded_at, temperature, humidity
FROM iot_sensor_readings
WHERE device_id = ? AND day = '2024-06-15'
  AND recorded_at > toTimestamp(now()) - 3600000;
      \end{lstlisting}
    \column{0.47\textwidth}
      \begin{tcolorbox}[deepdivebox, title={Deep Dive: Partition Size \& Bucketing}]
        \scriptsize
        \textbf{Problem}: a sensor emitting 10 readings/second for 1 year creates $\approx 315$ million cells in one partition — exceeding the recommended 100K cells/partition limit.
        \vspace{0.1cm}\\
        \textbf{Solution — Time Bucketing}: add a time component (\texttt{day}, \texttt{week}, or \texttt{month}) to the partition key. Each day's readings form a separate partition, keeping each partition at a manageable size ($\approx 864K$ cells/day for 10 Hz).
        \vspace{0.1cm}\\
        Queries spanning multiple days must issue one query per bucket and merge results in the application layer — a standard trade-off in Cassandra design.
      \end{tcolorbox}
      \begin{tcolorbox}[cassbox, title=TWCS Compaction]
        \scriptsize
        \textbf{TimeWindowCompactionStrategy} groups SSTables by time window and compacts only within the same window — perfectly aligned with time-bucketed TTL data. Avoids compacting expired data with live data.
      \end{tcolorbox}
  \end{columns}
\end{frame}

% ------------------------------------------------------------------
\begin{frame}{Cassandra Use Cases \& Anti-Patterns}
  \begin{columns}[T]
    \column{0.48\textwidth}
      \begin{tcolorbox}[cassbox, title={[OK] Ideal Use Cases}]
        \small
        \begin{itemize}
          \item \textbf{IoT sensor data} — time-bucketed, high-frequency writes from millions of devices with automatic TTL expiry
          \item \textbf{Chat / messaging systems} — partition by room ID + cluster by timestamp; sequential reads serve chat history instantly
          \item \textbf{User activity logs} — clickstreams, event logs, session tracking at massive scale
          \item \textbf{Leaderboards} — partition by game/category, cluster by score (DESC)
          \item \textbf{Inventory \& catalog} — wide rows accommodate product variant attributes without schema changes
          \item \textbf{Multi-region active-active} — NetworkTopologyStrategy + LOCAL\_QUORUM enables true active-active with no single region as master
        \end{itemize}
      \end{tcolorbox}
    \column{0.48\textwidth}
      \begin{tcolorbox}[deepdivebox, title={[!] Anti-Patterns}]
        \small
        \begin{itemize}
          \item \textbf{Relationship-heavy queries} — ``find all users who purchased X and also rated Y above 4 stars'' requires scanning multiple partitions or application-side filtering
          \item \textbf{Ad-hoc analytics} — unknown WHERE clause patterns require ALLOW FILTERING, which scans all partitions — catastrophically slow at scale
          \item \textbf{Frequent updates to the same cell} — creates a chain of versions that slow down reads (compaction hasn't merged them yet)
          \item \textbf{Many small partitions} — millions of single-row partitions defeat the purpose of the partition model; consider column-family restructuring
          \item \textbf{Counting} — \texttt{SELECT COUNT(*)} is a full table scan; use \texttt{counter} columns or external counters (Redis) instead
        \end{itemize}
      \end{tcolorbox}
  \end{columns}
\end{frame}

% ==================================================================
% SECTION 4: DEEP COMPARISON
% ==================================================================
\section{Deep Comparison}

\begin{frame}{Architectural Comparison: At a Glance}
  \scriptsize
  \renewcommand{\arraystretch}{1.5}
  \begin{tabular}{>{\bfseries}p{2.5cm} p{5.0cm} p{5.0cm}}
    \toprule
    \textbf{Dimension} & \textbf{\textcolor{neo4jGreen}{Neo4j}} & \textbf{\textcolor{cassBlue}{Apache Cassandra}} \\
    \midrule
    Category & Graph Database & Wide-Column Store \\
    Origin & Sweden, 2007 & Facebook, 2007 (Apache 2009) \\
    CAP position & CA / CP (Raft-based) & AP (tunable) \\
    Consistency model & ACID, full transactions & BASE, eventual (tunable to strong) \\
    Topology & Master-follower (Core + Read Replicas) & Masterless ring (peer-to-peer) \\
    Data model & Nodes, Relationships, Properties & Partitions, Rows, Clustering Columns \\
    Query language & Cypher (openCypher / GQL) & CQL (SQL-like) \\
    Write pattern & Raft-committed, ACID & Append to Commit Log + Memtable \\
    Read pattern & Index-free adjacency pointer chain & Bloom filter + Partition index + SSTable merge \\
    Scalability axis & Scale-up (+ Fabric for sharding) & Linear scale-out (add nodes to ring) \\
    Fault tolerance & Raft quorum (majority must survive) & RF copies; any node can serve any request \\
    Schema flexibility & High (schema-optional labels \& props) & Moderate (per-table schema, strict partition key) \\
    \bottomrule
  \end{tabular}
\end{frame}

% ------------------------------------------------------------------
\begin{frame}{Performance Characteristics: Write Path}
  \begin{columns}[T]
    \column{0.48\textwidth}
      \begin{tcolorbox}[neo4jbox, title=Neo4j Write Performance]
        \small
        \textbf{Strengths:}
        \begin{itemize}\scriptsize
          \item Transactional integrity — complex multi-node/multi-relationship writes in a single ACID transaction
          \item Relationship creation is $O(1)$: append two relationship records and update two node pointers
        \end{itemize}
        \textbf{Bottlenecks:}
        \begin{itemize}\scriptsize
          \item Raft consensus adds latency: $\sim$5--50ms extra per write to achieve quorum
          \item Pessimistic locking can cause contention on ``hot'' nodes (supernode problem)
          \item Bulk loading of millions of nodes requires offline batch import tools (\texttt{neo4j-admin import}) rather than transactional Cypher
        \end{itemize}
      \end{tcolorbox}
    \column{0.48\textwidth}
      \begin{tcolorbox}[cassbox, title=Cassandra Write Performance]
        \small
        \textbf{Strengths:}
        \begin{itemize}\scriptsize
          \item Writes are \textbf{sequential I/O} to Commit Log + Memtable — extremely fast ($<$1ms latency at ONE)
          \item No locking: last-write-wins via timestamp means concurrent writes never block
          \item Linear throughput scaling: doubling nodes roughly doubles write capacity
          \item Benchmark: single node can sustain \textbf{100,000+ writes/second}
        \end{itemize}
        \textbf{Bottlenecks:}
        \begin{itemize}\scriptsize
          \item Tombstone accumulation from frequent deletes/updates degrades read performance
          \item Compaction I/O spikes can temporarily consume disk bandwidth
        \end{itemize}
      \end{tcolorbox}
  \end{columns}
\end{frame}

% ------------------------------------------------------------------
\begin{frame}{Performance Characteristics: Read Path}
  \begin{columns}[T]
    \column{0.48\textwidth}
      \begin{tcolorbox}[neo4jbox, title=Neo4j Read Performance]
        \small
        \textbf{Strengths:}
        \begin{itemize}\scriptsize
          \item \textbf{Deep traversal} — 6-hop friendship query on a 1B node graph completes in milliseconds via pointer chasing
          \item Traversal time is proportional to \textit{result size}, not graph size
          \item In-memory page cache eliminates I/O for hot subgraphs
          \item Cypher planner optimizes complex pattern queries automatically
        \end{itemize}
        \textbf{Bottlenecks:}
        \begin{itemize}\scriptsize
          \item Analytical queries over the full graph (e.g., global PageRank) are memory-intensive; GDS library manages in-memory projections
          \item Read Replicas can lag behind Core by a configurable delta
        \end{itemize}
      \end{tcolorbox}
    \column{0.48\textwidth}
      \begin{tcolorbox}[cassbox, title=Cassandra Read Performance]
        \small
        \textbf{Strengths:}
        \begin{itemize}\scriptsize
          \item \textbf{Single partition reads} are extremely fast: Bloom filter eliminates most disk seeks; clustering columns pre-sort data for sequential reads
          \item Time-series queries with LIMIT on a clustered partition are $O(1)$ — they read from the head of the sorted SSTable
        \end{itemize}
        \textbf{Bottlenecks:}
        \begin{itemize}\scriptsize
          \item \textbf{Multi-partition reads} — queries crossing partitions require coordinating multiple nodes; latency scales with partition count
          \item \texttt{ALLOW FILTERING} causes a \textit{full table scan} — catastrophic at scale; always a design smell
          \item Read amplification increases with SSTable count (mitigated by compaction)
        \end{itemize}
      \end{tcolorbox}
  \end{columns}
\end{frame}

% ------------------------------------------------------------------
\begin{frame}{Scalability: Scale-Up vs.\ Scale-Out}
  \begin{columns}[T]
    \column{0.50\textwidth}
      \begin{tcolorbox}[neo4jbox, title=Neo4j: Vertical Scaling (Scale-Up)]
        \small
        \begin{itemize}
          \item Core strength: adding RAM to fit more of the graph in the \textbf{page cache} provides near-linear read performance gains
          \item Raft Core Servers must be \textit{identical in capability} — the weakest node determines consensus speed
          \item Horizontal write scaling via \textbf{Fabric} sharding, but requires application-level query routing and increased operational complexity
          \item Practical limit: a single Core Server can handle graphs of $\sim$\textbf{billions of nodes and relationships} if RAM is sufficient ($\sim$hundreds of GB)
        \end{itemize}
      \end{tcolorbox}
    \column{0.47\textwidth}
      \begin{tcolorbox}[cassbox, title=Cassandra: Linear Scale-Out]
        \small
        \begin{itemize}
          \item Adding a node to the ring \textbf{automatically redistributes} token ranges; capacity scales nearly linearly
          \item Both \textit{reads and writes} benefit from adding nodes — unlike primary/replica architectures where only reads scale
          \item Netflix demonstrated linear scaling from 48 to 288 nodes with near-linear throughput increase
          \item Operational simplicity: \texttt{nodetool bootstrap} / \texttt{decommission} handle node lifecycle without downtime
        \end{itemize}
      \end{tcolorbox}
      \vspace{0.2cm}
      \begin{tcolorbox}[deepdivebox, title=Key Insight]
        \scriptsize
        If your bottleneck is \textbf{relationship traversal depth}: scale up Neo4j RAM. If your bottleneck is \textbf{write throughput or storage volume}: add Cassandra nodes.
      \end{tcolorbox}
  \end{columns}
\end{frame}

% ------------------------------------------------------------------
\begin{frame}{Consistency \& Durability Trade-offs}
  \begin{columns}[T]
    \column{0.50\textwidth}
      \small
      \textbf{When strong consistency is non-negotiable:}
      \begin{itemize}
        \item Financial transactions (double-spend prevention)
        \item Access control decisions (security policy changes must be immediately visible)
        \item Inventory management (two users must not both think the last item is available)
      \end{itemize}
      \vspace{0.2cm}
      \textbf{Neo4j} handles these naturally via full ACID transactions.
      \vspace{0.2cm}
      \textbf{Cassandra} can handle these via:
      \begin{itemize}\scriptsize
        \item \texttt{QUORUM + QUORUM} (strong consistency) — sacrifices some availability
        \item \textbf{Lightweight Transactions} (Paxos) — full linearizability, but $4\times$ slower
      \end{itemize}
    \column{0.47\textwidth}
      \begin{tcolorbox}[deepdivebox, title={Deep Dive: Last-Write-Wins Conflicts}]
        \scriptsize
        Cassandra resolves concurrent writes using \textbf{timestamps}. The write with the higher timestamp wins. However:
        \begin{itemize}
          \item Clock skew across nodes (even with NTP) can cause the \textit{earlier} write to win if its node's clock is ahead
          \item Solution: use \textbf{client-generated timestamps} and ensure monotonically increasing values
          \item The \texttt{USING TIMESTAMP} clause allows explicit control
        \end{itemize}
        \vspace{0.1cm}
        For true conflict-free data structures, Cassandra 4.0+ introduced \textbf{CRDT-based counters}, but general-purpose CRDTs require application design.
      \end{tcolorbox}
  \end{columns}
\end{frame}

% ------------------------------------------------------------------
\begin{frame}{Operational Complexity Comparison}
  \scriptsize
  \renewcommand{\arraystretch}{1.5}
  \begin{tabular}{>{\bfseries}p{2.8cm} p{4.5cm} p{4.5cm}}
    \toprule
    \textbf{Operation} & \textbf{\textcolor{neo4jGreen}{Neo4j}} & \textbf{\textcolor{cassBlue}{Cassandra}} \\
    \midrule
    Initial setup & Simple single-instance; Causal Clustering adds complexity & Requires cluster planning (RF, DC layout, vnodes) upfront \\
    Schema changes & Add labels/properties without downtime; index creation is online & \texttt{ALTER TABLE} is instant (metadata only); data backfill is separate \\
    Backup & Online backup via \texttt{neo4j-admin backup}; consistent snapshot & Snapshots per-node via \texttt{nodetool snapshot}; coordination needed \\
    Monitoring & Built-in metrics endpoint, Prometheus integration, Neo4j Ops Manager & \texttt{nodetool} CLI, JMX metrics, DataStax Metrics Collector \\
    Upgrades & Rolling upgrade supported (Core + Read Replica) & Rolling upgrade: drain, upgrade, restart node-by-node; fully online \\
    Repair & Not required (Raft ensures consistency) & \texttt{nodetool repair} must run periodically to reconcile replicas \\
    Tuning levers & Page cache size, heap, query plan hints & Compaction strategy, memtable size, CL per query, GC grace \\
    Cloud-managed & AuraDB (fully managed, Neo4j Inc.) & Astra DB (DataStax), Amazon Keyspaces (AWS) \\
    \bottomrule
  \end{tabular}
\end{frame}

% ------------------------------------------------------------------
\begin{frame}{Developer Experience Comparison}
  \begin{columns}[T]
    \column{0.48\textwidth}
      \begin{tcolorbox}[neo4jbox, title=Neo4j Developer Experience]
        \small
        \begin{itemize}
          \item \textbf{Neo4j Browser / Bloom} — visual graph explorer; great for understanding data relationships intuitively
          \item \textbf{Cypher readability} — ``I want users who bought X and rated Y'' translates almost directly into a readable Cypher pattern
          \item \textbf{Schema flexibility} — add properties and labels to nodes without migration scripts
          \item \textbf{Driver ecosystem} — official drivers for Python, Java, JavaScript, Go, .NET; community drivers for many others
          \item \textbf{GraphQL integration} — Neo4j GraphQL Library auto-generates schema from graph model
        \end{itemize}
      \end{tcolorbox}
    \column{0.48\textwidth}
      \begin{tcolorbox}[cassbox, title=Cassandra Developer Experience]
        \small
        \begin{itemize}
          \item \textbf{CQL familiarity} — SQL-like syntax lowers onboarding barrier for developers with SQL background
          \item \textbf{Strict modeling discipline} — query-first design requires upfront query enumeration; changing queries later may require new tables
          \item \textbf{Driver ecosystem} — DataStax Java Driver, Python Driver, Node.js Driver; all support async, connection pooling, speculative execution
          \item \textbf{Cassandra Workload Analyzer} — tool for identifying anti-patterns in CQL queries
          \item \textbf{Spring Data Cassandra} — Spring Boot integration for rapid application development
        \end{itemize}
      \end{tcolorbox}
  \end{columns}
\end{frame}

% ------------------------------------------------------------------
\begin{frame}{Security Features Comparison}
  \begin{columns}[T]
    \column{0.48\textwidth}
      \begin{tcolorbox}[neo4jbox, title=Neo4j Security]
        \small
        \begin{itemize}
          \item \textbf{Authentication}: native username/password, LDAP, SSO (OIDC)
          \item \textbf{Authorization}: role-based access control (RBAC) with fine-grained privileges at database, graph, and property levels
          \item \textbf{Property-level security} (Enterprise): hide specific node/relationship properties from users without sufficient privilege
          \item \textbf{Subgraph-level security}: restrict which nodes/relationships a role can see
          \item \textbf{Encryption}: TLS for bolt/HTTP protocols; data-at-rest encryption via OS/storage layer
          \item \textbf{Audit logging}: full query audit log for compliance (GDPR, HIPAA)
        \end{itemize}
      \end{tcolorbox}
    \column{0.48\textwidth}
      \begin{tcolorbox}[cassbox, title=Cassandra Security]
        \small
        \begin{itemize}
          \item \textbf{Authentication}: internal (CassandraAutenticator) or LDAP via DSE
          \item \textbf{Authorization}: object-level permissions (keyspace, table, function) via \texttt{GRANT/REVOKE}
          \item \textbf{SSL/TLS}: node-to-node and client-to-node encryption configurable independently
          \item \textbf{Transparent Data Encryption (TDE)}: DSE Enterprise; encrypts SSTable files on disk
          \item \textbf{Audit logging}: introduced in Cassandra 4.0; logs CQL operations per-keyspace
          \item \textbf{Network}: firewall rules recommended; no native row-level security without DSE
        \end{itemize}
      \end{tcolorbox}
  \end{columns}
\end{frame}

% ------------------------------------------------------------------
\begin{frame}{Ecosystem \& Integrations}
  \begin{columns}[T]
    \column{0.50\textwidth}
      \begin{tcolorbox}[neo4jbox, title=Neo4j Ecosystem]
        \scriptsize
        \begin{itemize}
          \item \textbf{GDS Library}: 65+ graph algorithms (centrality, community, pathfinding, ML embeddings)
          \item \textbf{APOC}: 450+ utility procedures
          \item \textbf{GraphQL API}: auto-generated from graph schema
          \item \textbf{Data connectors}: Kafka, Spark (GraphFrames), Elasticsearch
          \item \textbf{ETL tools}: Neo4j ETL tool for RDBMS migration; Data Importer (GUI)
          \item \textbf{Visualization}: Neo4j Bloom, yFiles, Graphistry (GPU-accelerated)
          \item \textbf{LLM integration}: LangChain Neo4j integration; Knowledge Graph + RAG patterns
        \end{itemize}
      \end{tcolorbox}
    \column{0.47\textwidth}
      \begin{tcolorbox}[cassbox, title=Cassandra Ecosystem]
        \scriptsize
        \begin{itemize}
          \item \textbf{Apache Spark (Cassandra Connector)}: run Spark analytics directly over Cassandra data
          \item \textbf{Apache Kafka (Kafka Connect)}: stream events into Cassandra at low latency
          \item \textbf{Presto/Trino}: SQL analytics layer over Cassandra for ad-hoc queries
          \item \textbf{Elasticsearch}: common pattern — Cassandra for primary store, ES for full-text search
          \item \textbf{Stargate}: API gateway providing REST, GraphQL, and Document APIs over Cassandra
          \item \textbf{Pulsar}: Apache Pulsar native Cassandra sink/source connector
          \item \textbf{Flink}: streaming analytics with Cassandra state store
        \end{itemize}
      \end{tcolorbox}
  \end{columns}
\end{frame}

% ==================================================================
% SECTION 5: POLYGLOT PERSISTENCE
% ==================================================================
\section{Polyglot Persistence}

\begin{frame}{Polyglot Persistence: The Modern Architecture Paradigm}
  \begin{columns}[T]
    \column{0.55\textwidth}
      \small
      \textbf{Definition:}
      \begin{quote}
        \textit{Polyglot persistence} is the practice of using different data storage technologies within a single application to match each component's specific access patterns and data shape requirements.
      \end{quote}
      \vspace{0.2cm}
      \textbf{Origin:} Coined by Martin Fowler and Pramod Sadalage in \textit{NoSQL Distilled} (2012).
      \vspace{0.3cm}

      \textbf{Core principle:} Abandon the idea of a single ``one-size-fits-all'' database. Instead:
      \begin{itemize}\small
        \item Identify the nature of each data domain
        \item Select the storage engine optimized for that domain's access pattern
        \item Accept operational overhead of managing multiple systems in exchange for each system performing at its theoretical optimum
      \end{itemize}
    \column{0.42\textwidth}
      \begin{tcolorbox}[infobox, title=Common Polyglot Stack]
        \scriptsize
        \begin{tabular}{@{}ll@{}}
          \textbf{Relational} & Core transactions (orders) \\
          \textbf{Redis} & Sessions, caching, pub/sub \\
          \textbf{Elasticsearch} & Full-text search \\
          \textbf{Neo4j} & Recommendation, fraud \\
          \textbf{Cassandra} & Events, logs, timeseries \\
          \textbf{S3 / Blob} & Binary assets \\
        \end{tabular}
      \end{tcolorbox}
      \vspace{0.2cm}
      \begin{alertblock}{Trade-off}
        \small Polyglot persistence increases \textbf{data consistency complexity} (cross-store transactions are not atomic) and \textbf{operational burden} (more systems to monitor and maintain). Microservices architectures absorb this naturally via service boundaries.
      \end{alertblock}
  \end{columns}
\end{frame}

% ------------------------------------------------------------------
\begin{frame}{Combining Neo4j and Cassandra: Architecture Patterns}
  \begin{columns}[T]
    \column{0.55\textwidth}
      \small
      \textbf{Pattern 1: Event Sourcing + Graph Projection}
      \begin{itemize}\scriptsize
        \item All user events (clicks, purchases, ratings) are written to \textbf{Cassandra} as immutable event log (partition by user\_id, cluster by event\_time)
        \item An asynchronous consumer reads the Cassandra changelog and projects relevant relationship events into \textbf{Neo4j} (e.g., ``User A purchased Product B'' → creates \texttt{[:PURCHASED]} edge)
        \item Neo4j serves real-time graph traversal queries (recommendations, fraud detection) over the projected graph
      \end{itemize}
      \vspace{0.2cm}
      \textbf{Pattern 2: Write to Both, Query from Right}
      \begin{itemize}\scriptsize
        \item Application writes to both stores in the same business transaction (eventually consistent across stores)
        \item Relationship queries → Neo4j; timeseries/bulk queries → Cassandra
        \item A Kafka topic acts as the bridge; consumers update each store independently
      \end{itemize}
    \column{0.42\textwidth}
      \begin{tcolorbox}[deepdivebox, title={Deep Dive: Data Consistency Risk}]
        \scriptsize
        Writing to two databases in the same business operation creates a \textbf{distributed transaction problem}. If the Cassandra write succeeds but the Neo4j write fails:
        \begin{itemize}
          \item The event log is correct (Cassandra)
          \item But the recommendation graph is missing an edge (Neo4j)
        \end{itemize}
        \textbf{Mitigation strategies:}
        \begin{enumerate}
          \item \textbf{Outbox pattern}: write to Cassandra first; publish event to Kafka; Neo4j consumer processes at-least-once
          \item \textbf{Saga pattern}: implement compensating transactions for rollback
          \item \textbf{Eventual consistency SLA}: accept a brief window where Neo4j lags Cassandra
        \end{enumerate}
      \end{tcolorbox}
  \end{columns}
\end{frame}

% ------------------------------------------------------------------
\begin{frame}{Case Study: University Social Platform Architecture}
  \begin{columns}[T]
    \column{0.55\textwidth}
      \small
      \textbf{``BK Connect'' — a university social network:}
      \vspace{0.2cm}

      \textbf{Service 1: Recommendation Engine (Neo4j)}
      \begin{itemize}\scriptsize
        \item Data: Student nodes, Course nodes, Faculty nodes
        \item Relationships: \texttt{[:ENROLLED\_IN]}, \texttt{[:FRIEND]}, \texttt{[:STUDIES\_WITH]}, \texttt{[:SHARES\_FACULTY]}
        \item Queries: ``Who should I connect with?'', ``What courses are popular in my network?''
        \item Latency requirement: $<$50ms for 3-hop traversals
      \end{itemize}
      \vspace{0.2cm}
      \textbf{Service 2: Messaging Service (Cassandra)}
      \begin{itemize}\scriptsize
        \item Table: \texttt{messages\_by\_room} (partition: room\_id, cluster: sent\_at DESC)
        \item Write pattern: 10,000+ messages/second during exam periods
        \item Query: ``Get last 50 messages in room X'' — answered by a single partition sequential read
        \item HA: RF=3 across two datacenters; LOCAL\_QUORUM guarantees zero downtime during a full DC failure
      \end{itemize}
    \column{0.42\textwidth}
      \textbf{Service 3: Activity Feed (Cassandra)}
      \begin{itemize}\scriptsize
        \item Denormalized per-user timeline table
        \item On each social action (post, like, comment), the event is fanned out to followers' timelines in Cassandra
        \item Partition: (user\_id, week\_bucket); Cluster: activity\_at DESC
      \end{itemize}
      \vspace{0.2cm}
      \textbf{Service 4: User Profiles (PostgreSQL)}
      \begin{itemize}\scriptsize
        \item Core identity data with full ACID — Neo4j and Cassandra reference user\_id as a foreign key concept
      \end{itemize}
      \vspace{0.2cm}
      \begin{tcolorbox}[infobox, title=Integration Fabric]
        \scriptsize
        \textbf{Apache Kafka} bridges all services: profile updates propagate to Neo4j and Cassandra via event streams; no service directly writes to another service's database.
      \end{tcolorbox}
  \end{columns}
\end{frame}

% ------------------------------------------------------------------
\begin{frame}{Decision Framework: Choosing the Right Database}
  \begin{columns}[T]
    \column{0.50\textwidth}
      \small
      \textbf{Choose Neo4j when:}
      \begin{itemize}\scriptsize
        \item Your primary questions are about \textit{relationships} (``how are X and Y connected?'', ``what is the shortest path?'')
        \item The graph structure is unknown at schema design time (exploratory graph analytics)
        \item You need ACID transactions over multi-entity graph operations
        \item Traversal depth matters more than raw throughput
        \item The domain is naturally graph-shaped: social networks, knowledge graphs, dependency trees, org hierarchies
      \end{itemize}
      \vspace{0.2cm}
      \textbf{Choose Cassandra when:}
      \begin{itemize}\scriptsize
        \item Write throughput is the primary concern ($>$10K writes/second)
        \item Data access patterns are known and stable
        \item Downtime is unacceptable (global 24/7 availability)
        \item Data volumes require linear horizontal scaling to petabytes
        \item Time-series, events, or log-style data is the primary payload
      \end{itemize}
    \column{0.47\textwidth}
      \begin{tcolorbox}[infobox, title=Decision Checklist]
        \scriptsize
        \begin{tabular}{@{}lcc@{}}
          & \textbf{Neo4j} & \textbf{Cassandra} \\
          \midrule
          Relationship queries? & \checkmark\checkmark & \texttimes \\
          High write throughput? & \texttimes & \checkmark\checkmark \\
          ACID transactions needed? & \checkmark\checkmark & (\checkmark) \\
          Graph traversal queries? & \checkmark\checkmark & \texttimes \\
          Time-series data? & \texttimes & \checkmark\checkmark \\
          Flexible queries? & \checkmark\checkmark & \texttimes \\
          Linear scale-out? & \texttimes & \checkmark\checkmark \\
          Multi-DC active-active? & (\checkmark) & \checkmark\checkmark \\
          Schema flexibility? & \checkmark\checkmark & \checkmark \\
          Known access patterns? & \checkmark & \checkmark\checkmark \\
        \end{tabular}
      \end{tcolorbox}
  \end{columns}
\end{frame}

% ==================================================================
% SECTION 6: EMERGING TRENDS
% ==================================================================
\section{Emerging Trends \& Research Directions}

\begin{frame}{Emerging Trends in Graph Databases}
  \begin{columns}[T]
    \column{0.50\textwidth}
      \begin{tcolorbox}[neo4jbox, title=Graph + AI / LLMs]
        \small
        \begin{itemize}
          \item \textbf{Knowledge Graph + RAG (Retrieval-Augmented Generation)}: LLMs are augmented by Neo4j knowledge graphs that provide structured, relationship-aware context — reducing hallucinations
          \item \textbf{GraphRAG} (Microsoft Research, 2024): community detection on knowledge graphs provides hierarchical document summaries for LLM context
          \item \textbf{Graph Neural Networks}: Neo4j GDS provides GraphSAGE and FastRP embeddings that can be used as input features for ML classifiers
        \end{itemize}
      \end{tcolorbox}
    \column{0.47\textwidth}
      \begin{tcolorbox}[neo4jbox, title=GQL: The ISO Graph Query Standard]
        \small
        \begin{itemize}
          \item \textbf{GQL (ISO/IEC 39075)} was ratified in April 2024 — the first international standard for property graph query languages
          \item Heavily influenced by Cypher (openCypher), SQL/PGQ (SQL:2023), and PGQL
          \item Neo4j 5.x roadmap includes GQL compliance
          \item Signals a maturation of the graph database ecosystem comparable to what SQL did for relational databases in the 1980s
        \end{itemize}
      \end{tcolorbox}
  \end{columns}
\end{frame}

% ------------------------------------------------------------------
\begin{frame}{Emerging Trends in Wide-Column Stores}
  \begin{columns}[T]
    \column{0.50\textwidth}
      \begin{tcolorbox}[cassbox, title=Cassandra 5.0: Storage-Attached Indexing (SAI)]
        \small
        \begin{itemize}
          \item SAI (Storage-Attached Indexing) replaces the legacy SASI and secondary indexes with a new architecture that co-locates index data with SSTable files
          \item Enables efficient \textbf{non-partition-key filtering} — previously the Achilles heel of Cassandra data modeling
          \item Supports vector search (Ann indexes) — enabling Cassandra to serve as a \textbf{vector database} alongside its traditional workloads
          \item Opens Cassandra to similarity search for ML feature stores and embedding retrieval
        \end{itemize}
      \end{tcolorbox}
    \column{0.47\textwidth}
      \begin{tcolorbox}[cassbox, title=Cassandra \& the Serverless Future]
        \small
        \begin{itemize}
          \item \textbf{DataStax Astra DB}: fully managed Cassandra with serverless billing (per-request, not per-node) — removes the operational burden of ring management
          \item \textbf{Amazon Keyspaces}: Cassandra-compatible managed service on AWS; separates compute and storage (S3-based)
          \item These serverless models decouple Cassandra's scalability benefits from its operational complexity — making it accessible to teams without dedicated DBA expertise
        \end{itemize}
      \end{tcolorbox}
  \end{columns}
\end{frame}

% ------------------------------------------------------------------
\begin{frame}{Research Frontiers: Graph Computing at Scale}
  \begin{columns}[T]
    \column{0.52\textwidth}
      \small
      \textbf{Open research problems:}
      \begin{itemize}
        \item \textbf{Distributed graph traversal}: efficiently traversing graphs partitioned across hundreds of nodes without expensive cross-partition hops (the ``edge-cut'' problem)
        \item \textbf{Dynamic graph algorithms}: maintaining PageRank, community detection results incrementally as edges are added/removed — without recomputing from scratch
        \item \textbf{Temporal graphs}: efficiently querying ``what did the graph look like at time T?'' — requires bitemporal graph models
        \item \textbf{Heterogeneous graph learning}: Graph Neural Networks that handle multiple node/edge types (as in property graphs) without collapsing to a homogeneous representation
        \item \textbf{Privacy-preserving graph analytics}: federated graph learning across organizations without sharing raw graph topology
      \end{itemize}
    \column{0.45\textwidth}
      \begin{tcolorbox}[deepdivebox, title={Research Context: GNN on Property Graphs}]
        \scriptsize
        Traditional GNNs (GraphSAGE, GCN) assume homogeneous graphs (one node type, one edge type). Real-world property graphs like Neo4j instances are \textbf{heterogeneous}: a node can be \texttt{:Student}, \texttt{:Course}, or \texttt{:Faculty}; an edge can be \texttt{[:ENROLLED\_IN]} or \texttt{[:FRIEND]}.
        \vspace{0.1cm}\\
        \textbf{Heterogeneous GNNs} (HAN, HGT, R-GCN) handle multiple types but require graph-type-specific aggregation functions. Active research area with significant implications for fraud detection and biomedical knowledge graphs.
      \end{tcolorbox}
  \end{columns}
\end{frame}

% ==================================================================
% SECTION 7: SUMMARY
% ==================================================================
\section{Summary \& Conclusions}

\begin{frame}{Best Practices: Neo4j}
  \begin{columns}[T]
    \column{0.50\textwidth}
      \begin{tcolorbox}[neo4jbox, title={[OK] Do}]
        \small
        \begin{itemize}
          \item \textbf{Model the domain as-is}: let entities become nodes and business relationships become typed edges with properties
          \item \textbf{Use indexes at entry points}: always index the properties you use in the first \texttt{MATCH} clause
          \item \textbf{Leverage MERGE for idempotent loading}: prevents duplicate data during re-imports
          \item \textbf{Size the page cache first}: aim to hold the full node and relationship stores in RAM for best traversal performance
          \item \textbf{Use \texttt{PROFILE}}: always inspect the query plan for unexpected NodeByLabelScan operations
          \item \textbf{Batch large writes}: use \texttt{CALL \{\} IN TRANSACTIONS} for bulk operations to avoid single giant transactions that lock the graph
        \end{itemize}
      \end{tcolorbox}
    \column{0.47\textwidth}
      \begin{tcolorbox}[deepdivebox, title={[!] Don't}]
        \small
        \begin{itemize}
          \item \textbf{Don't use Neo4j as a logging store}: it is not designed for sequential, flat data ingestion; every write creates relationship pointers — expensive at billions of events
          \item \textbf{Don't create supernodes carelessly}: a node with millions of relationships degrades traversal. Partition relationships by type, date, or status
          \item \textbf{Don't skip relationship types}: an untyped \texttt{-[r]->} relationship loses semantic value and forces runtime filtering
          \item \textbf{Don't ignore the Raft quorum requirement}: in a 3-node cluster, 2 nodes must be healthy for writes. Plan accordingly for maintenance windows
        \end{itemize}
      \end{tcolorbox}
  \end{columns}
\end{frame}

% ------------------------------------------------------------------
\begin{frame}{Best Practices: Cassandra}
  \begin{columns}[T]
    \column{0.50\textwidth}
      \begin{tcolorbox}[cassbox, title={[OK] Do}]
        \small
        \begin{itemize}
          \item \textbf{Design tables for queries}: list all read queries first; design one table per query pattern
          \item \textbf{Choose partition keys that distribute evenly}: avoid partition keys that concentrate data on a few nodes (``hot spots'')
          \item \textbf{Use TTL for ephemeral data}: let Cassandra automatically expire sessions, caches, and rate-limit counters rather than writing explicit deletes (which become tombstones)
          \item \textbf{Run \texttt{nodetool repair} regularly}: keeps replicas in sync; frequency depends on \texttt{gc\_grace\_seconds} setting
          \item \textbf{Use \texttt{LOCAL\_QUORUM} for multi-DC}: balances strong consistency within a DC with low cross-DC latency
        \end{itemize}
      \end{tcolorbox}
    \column{0.47\textwidth}
      \begin{tcolorbox}[deepdivebox, title={[!] Don't}]
        \small
        \begin{itemize}
          \item \textbf{Don't use \texttt{ALLOW FILTERING} in production}: it triggers a full table scan that bypasses all partitioning logic — catastrophic at scale
          \item \textbf{Don't generate excessive tombstones}: frequent overwrites of nullable columns create read-amplifying tombstone chains; prefer \texttt{TTL} over explicit deletes
          \item \textbf{Don't ignore partition size}: aim for partitions between 1MB and 100MB. Very large partitions cause GC pressure and compaction stalls
          \item \textbf{Don't rely on \texttt{SELECT COUNT(*)}}: full-table count is a distributed scan; use approximate counters or counter tables instead
        \end{itemize}
      \end{tcolorbox}
  \end{columns}
\end{frame}

% ------------------------------------------------------------------
\begin{frame}{Summary: Core Takeaways}
  \begin{columns}[T]
    \column{0.50\textwidth}
      \begin{tcolorbox}[neo4jbox, title=Neo4j in One Paragraph]
        \small
        Neo4j is the tool you reach for when \textbf{relationships are the data}. Its property graph model and index-free adjacency make multi-hop traversals practically instant, independent of total graph size. Full ACID transactions ensure graph integrity. Cypher's ASCII-art syntax makes complex relationship queries readable. The cost: write throughput and horizontal scalability are more constrained than purpose-built distributed stores.
      \end{tcolorbox}
      \vspace{0.3cm}
      \begin{tcolorbox}[cassbox, title=Cassandra in One Paragraph]
        \small
        Cassandra is the tool you reach for when \textbf{volume, velocity, and uptime are non-negotiable}. Its masterless ring, tunable consistency, and append-only write path deliver linear scalability and zero single points of failure. The cost: you must model your tables around known query patterns, accept denormalization, and abandon joins.
      \end{tcolorbox}
    \column{0.47\textwidth}
      \begin{tcolorbox}[infobox, title=The Design Philosophy Difference]
        \normalsize
        \begin{center}
          \textbf{Neo4j:}\\
          \textit{``Model the reality\\of your domain.''}
        \end{center}
        \vspace{0.3cm}
        \begin{center}
          \textbf{Cassandra:}\\
          \textit{``Model the queries\\your application needs.''}
        \end{center}
        \vspace{0.3cm}
        \hrule
        \vspace{0.3cm}
        Both are correct philosophies — for \textit{different} classes of problems. Wisdom is knowing which class your problem belongs to.
      \end{tcolorbox}
  \end{columns}
\end{frame}

% ------------------------------------------------------------------
\begin{frame}{Conclusion: The Bigger Picture}
  \begin{columns}[T]
    \column{0.60\textwidth}
      \small
      \textbf{What this course should leave you with:}
      \begin{itemize}
        \item There is \textbf{no universally best database}. Every system makes trade-offs that make it optimal for some workloads and terrible for others.
        \item The CAP theorem, ACID vs.\ BASE, and the relational vs.\ NoSQL divide are not just academic distinctions — they are \textbf{engineering decisions with real production consequences}.
        \item \textbf{Polyglot persistence} is the pragmatic synthesis: combine Neo4j's graph intelligence with Cassandra's distributed power, bridged by event streaming, to build systems greater than the sum of their parts.
        \item The emerging intersection of \textbf{graph databases + LLMs} (GraphRAG, Knowledge Graphs) is a significant research and industry frontier — highly relevant if you are pursuing graduate research in AI or distributed systems.
        \item Always \textbf{benchmark with your actual workload}. Theoretical characteristics are starting points; production performance is empirical.
      \end{itemize}
    \column{0.38\textwidth}
      \begin{tcolorbox}[infobox, title=Further Reading]
        \scriptsize
        \begin{itemize}
          \item \textit{Graph Databases}, Robinson, Webber \& Eifrem (O'Reilly)
          \item \textit{NoSQL Distilled}, Fowler \& Sadalage (Addison-Wesley)
          \item \textit{Cassandra: The Definitive Guide}, Carpenter \& Hewitt (O'Reilly)
          \item Brewer (2000): ``Towards Robust Distributed Systems'' (CAP theorem)
          \item Abadi (2012): ``Consistency Tradeoffs in Modern Distributed Databases'' (PACELC)
          \item Edge et al.\ (2024): ``GraphRAG: Unlocking LLM Discovery on Narrative Private Data'' (Microsoft)
          \item ISO/IEC 39075:2024 (GQL Standard)
        \end{itemize}
      \end{tcolorbox}
  \end{columns}
\end{frame}

% ------------------------------------------------------------------
\begin{frame}{}
  \begin{center}
    \vspace{1.5cm}
    {\Huge \textbf{Thank You}}\\[0.5cm]
    {\large Questions \& Discussion}\\[1.0cm]
    \hrule
    \vspace{0.5cm}
    \begin{columns}[c]
      \column{0.45\textwidth}
        \begin{tcolorbox}[neo4jbox]
          \centering\small
          \textbf{Neo4j}\\
          \textit{``If your data is a graph,\\model it as a graph.''}
        \end{tcolorbox}
      \column{0.45\textwidth}
        \begin{tcolorbox}[cassbox]
          \centering\small
          \textbf{Cassandra}\\
          \textit{``Scale without\\compromise on uptime.''}
        \end{tcolorbox}
    \end{columns}
    \vspace{0.5cm}
    {\small L\^{a}m Kh\'{a}nh Tr\`{i}nh $\cdot$ HCMUT $\cdot$ \texttt{CO3021 — Database Systems}}
  \end{center}
\end{frame}

\end{document}